\section*{الفصل \ref{ch:b1:field-lab-methods-1} --- طرائق الحقل والمختبر I}

\begin{solution}{exo:b1:field-lab-methods-1:1}
عند 1:10\,000 يصبح \qty{4.9}{m} قدره $4.9/10\,000 = \qty{0.49}{mm}$؛ وعند 1:25\,000 يكون $4.9/25\,000 = \qty{0.20}{mm}$، أي نحو عرض خط القلم الرصاص.
\end{solution}

\begin{solution}{exo:b1:field-lab-methods-1:2}
عند ريكيافيك $D = \ang{-11.28}$: $52 - 11.28 = \ang{040.7}$. وعند ويلينغتون $D = \ang{+23.23}$: $52 + 23.23 = \ang{075.2}$. فالقراءة البوصلية نفسها تعني اتجاهين يبعدان \ang{34.5} عن بعضهما.
\end{solution}

\begin{solution}{exo:b1:field-lab-methods-1:3}
$20\to015$. ومن أجل $t = \ang{15}$ و$p = \ang{20}$، تكون متجهة الوحدة الهابطة
\[
  (\cos\ang{20}\sin\ang{15},\ \cos\ang{20}\cos\ang{15},\ -\sin\ang{20}) =
  (0.2432,\ 0.9077,\ -0.3420).
\]
\end{solution}

\begin{solution}{exo:b1:field-lab-methods-1:4}
المتوسط الحسابي هو $(350 + 358 + 6)/3 = \ang{238}$، أي شبه معاكس لكل القراءات الثلاث. وتُجمع المتجهات الأفقية الوحدوية $(\sin a, \cos a)$ فتعطي
$(-0.1736 - 0.0349 + 0.1045,\ 0.9848 + 0.9994 + 0.9945) = (-0.1040, 2.9787)$،
وسمتها هو $\operatorname{atan2}(-0.1040, 2.9787) = \ang{-2.0}$، أي \ang{358.0}: هو المتوسط الصحيح، في وسط القراءات.
\end{solution}

\begin{solution}{exo:b1:field-lab-methods-1:5}
المتوسط $175/5 = \ang{35.0}$. الانحرافات $-1, 2, 0, -2, 1$؛ ومجموع المربعات 10؛ و$s = \sqrt{10/4} = \ang{1.58}$. والخطأ المعياري $1.58/\sqrt5 = \ang{0.71}$.
\end{solution}

\begin{solution}{exo:b1:field-lab-methods-1:6}
النواظم هي $(\sin\ang{30}, 0, \cos\ang{30}) = (0.5, 0, 0.866)$ و$(-0.5, 0, 0.866)$؛ و$\vec S = (0, 0, 1.732)$ شاقولية: فالمتوسط مستوٍ أفقي، و$\bar R = 1.732/2 = 0.866$، وكل قراءة تبعد \ang{30} عن المتوسط. والقراءتان لا تأخذان عيّنة من سطح واحد بأخطاء: فهما من جناحين \omterm{def:b1:geological-maps:fold}{لطية}، أي مجتمعان مختلفان، ومتوسطهما لا يصف أيًّا منهما. فيُتوسَّط كل جناح على حدة؛ ومعًا يحدّدان \omterm{def:b1:geological-maps:fold}{محور الطية}، وهو الخط المشترك للمستويين، وهو هنا أفقي شمالي-جنوبي.
\end{solution}

\begin{solution}{exo:b1:field-lab-methods-1:7}
$\tan p = \tan\ang{12}\cos(\ang{80} - \ang{120}) = 0.2126\times0.7660 =
0.1628$: $p = \ang{9.2}$، فالخدوش تغوص نحو \ang{120}. أما من أجل الاتجاه \ang{170}، فإن $\cos(\ang{-90}) = 0$: فالخط أفقي، على طول \omterm{def:b1:geological-maps:attitude}{الامتداد}.
\end{solution}

\begin{solution}{exo:b1:field-lab-methods-1:8}
$0.47/0.20 = 2.35$ سنة بعد 1 يناير 2026، أي في ربيع 2028. وفي 1 يناير 2030: $-0.47 + 4\times0.20 = \ang{+0.33}$، شرق الشمال الجغرافي.
\end{solution}

\begin{solution}{exo:b1:field-lab-methods-1:9}
تقع النواظم الثلاثة في المستوي الشاقولي ذي \omterm{def:b1:spherical-astronomy:horizontal}{السمت} \ang{080}، على بُعد \ang{10} و\ang{12} و\ang{14} عن الشاقول؛ والطرفان متماثلان حول الوسط، فمجموعهما يقع عليه: والمتوسط هو $080/12$، مع $\theta = \ang{2}, 0, \ang{2}$. ثم $\bar R = (1 +
2\cos\ang{2})/3 = (1 + 1.998782)/3 = 0.999594$، و$\sqrt{2(1 - \bar R)} =
\sqrt{0.000812} = \qty{0.02850}{rad} = \ang{1.633}$؛ وبدقة، $\theta_{\text{جذر المتوسط}} = \sqrt{(4 + 0 + 4)/3} = \ang{1.633}$.
\end{solution}

\begin{solution}{exo:b1:field-lab-methods-1:10}
المقطع دائرة نصف قطرها $\sqrt{r^2 - h^2}$. ومتوسطها
$\frac1r\int_0^r\sqrt{r^2 - h^2}\,\mathrm dh = \frac1r\cdot\frac{\pi r^2}{4}
= \frac{\pi r}{4}$، لأن التكامل هو مساحة ربع قرص. وحبيبات عرضها \qty{0.50}{mm} تبدو في المتوسط بعرض $0.50\times\pi/4 = \qty{0.39}{mm}$: \omterm{def:b1:field-lab-methods-1:thin-section}{فالشريحة الرقيقة} تقلّل من تقدير أحجام الحبيبات.
\end{solution}

\begin{solution}{exo:b1:field-lab-methods-1:11}
الدوران بالزاوية $D$ حول الشاقول،
$(x, y, z) \mapsto (x\cos D + y\sin D,\ -x\sin D + y\cos D,\ z)$، يحمل
$(\sin\delta\sin\alpha, \sin\delta\cos\alpha, \cos\delta)$ إلى
$(\sin\delta\sin(\alpha + D), \sin\delta\cos(\alpha + D), \cos\delta)$: فهو يضيف $D$ إلى كل اتجاه انحدار. ولأنه خطّي، يحمل المحصلة $\vec S$ إلى محصلة النواظم المدوّرة؛ وهو يحفظ الأطوال والمركبة الشاقولية، \omterm{def:b1:geological-maps:attitude}{فالانحدار} المتوسط $\arccos(S_U/\|\vec S\|)$ لا يتغيّر، ويدير الجزء الأفقي من $\vec S$ بالزاوية $D$، فيزداد \omterm{def:b1:geological-maps:attitude}{اتجاه الانحدار} المتوسط بمقدار $D$. وتصحيح كل قراءة على حدة بالانحراف، أو تصحيح متوسطها وحده، يعطي النتيجة نفسها.
\end{solution}

\begin{solution}{exo:b1:field-lab-methods-1:12}
التشتت على طول خط \omterm{def:b1:geological-maps:attitude}{الانحدار} الأشدّ هو تشتت في \omterm{def:b1:geological-maps:attitude}{الانحدار}: فالخطأ المعياري $2/\sqrt N \le 0.5$ من أجل $N \ge 16$. وعبره، يُدير دوران الناظم بمقدار \ang{2} \omterm{def:b1:geological-maps:attitude}{اتجاه الانحدار} بمقدار $2/\sin\ang{12} = \ang{9.62}$؛ والخطأ المعياري لمتوسط \omterm{def:b1:geological-maps:attitude}{اتجاه الانحدار} هو $9.62/\sqrt N \le 1$ من أجل $N \ge 92.5$، أي 93 قراءة. \omterm{def:b1:geological-maps:attitude}{فالانحدار} اللطيف يُعرف بسرعة؛ أما اتجاهه فلا.
\end{solution}

\begin{solution}{pb:b1:field-lab-methods-1:1}
\textbf{1.} اتجاهات \omterm{def:b1:geological-maps:attitude}{الانحدار}: $967/12 = \ang{80.6}$؛ الانحدارات: $145/12 =
\ang{12.1}$.

\textbf{2.} المتوسط الحسابي ل\ang{355} و\ang{005} هو \ang{180}، أي الاتجاه المعاكس. وهنا تقع كل اتجاهات \omterm{def:b1:geological-maps:attitude}{الانحدار} بين \ang{076} و\ang{085}، بعيدًا عن القفزة عند \ang{000}، فالمتوسط الحسابي يصلح.

\textbf{3.} نتحقّق من أن المتجهة الناظمية ذات طول واحد: $\|\vec n\|^2 = \sin^2\delta(\sin^2\alpha + \cos^2\alpha) +
\cos^2\delta = 1$. وخط \omterm{def:b1:geological-maps:attitude}{الانحدار} الهابط له الاتجاه
$(\cos\delta\sin\alpha, \cos\delta\cos\alpha, -\sin\delta)$، وجداؤه السلّمي مع $\vec n$ هو $\sin\delta\cos\delta - \cos\delta\sin\delta = 0$؛ وخط \omterm{def:b1:geological-maps:attitude}{الامتداد} أفقي وعمودي على \omterm{def:b1:geological-maps:attitude}{اتجاه الانحدار}، $(\cos\alpha, -\sin\alpha, 0)$، وجداؤه السلّمي مع $\vec n$ هو
$\sin\delta\sin\alpha\cos\alpha - \sin\delta\cos\alpha\sin\alpha = 0$.

\textbf{4.} القراءة 11 هي القراءة 2 نفسها: $(0.2222, 0.0352, 0.9744)$.
والقراءة 12:
\[
  (\sin\ang{11}\sin\ang{82},\ \sin\ang{11}\cos\ang{82},\ \cos\ang{11}) =
  (0.1890,\ 0.0266,\ 0.9816).
\]

\textbf{5.} $\vec S = (2.4747,\ 0.4082,\ 11.7294)$.

\textbf{6.} $\|\vec S\| = \sqrt{6.1241 + 0.1666 + 137.5788} = 11.99457$، و
$\bar R = 11.99457/12 = 0.999547$.

\textbf{7.} $\bar\alpha = \operatorname{atan2}(2.4747, 0.4082) = \ang{80.6}$
و$\bar\delta = \arccos(11.7294/11.99457) = \arccos 0.977893 = \ang{12.1}$:
فالمتوسط المتجهي هو $080.6/12.1$، وهو نفسه المتوسط الحسابي هنا.

\textbf{8.} اتجاهات \omterm{def:b1:geological-maps:attitude}{الانحدار}: $s = \ang{2.94}$؛ الانحدارات: $s = \ang{1.68}$.

\textbf{9.} $\theta_{\text{جذر المتوسط}} \approx \sqrt{2\times0.000453} =
\qty{0.0301}{rad} = \ang{1.72}$.

\textbf{10.} $\sqrt{1.68^2 + (\sin\ang{12.1}\times2.94)^2} = \sqrt{2.82 +
0.38} = \ang{1.79}$، قريب من \ang{1.72}: فالقراءات تسلك كما يقول نموذج الخطأ، ضمن أثر التقريب إلى الدرجة وعيّنة صغيرة.

\textbf{11.} \omterm{def:b1:geological-maps:attitude}{الانحدار}: يسهم ب2.82 من أصل 3.20 (بالدرجات المربعة). ويدخل \omterm{def:b1:geological-maps:attitude}{اتجاه الانحدار} من خلال $\sin\delta = 0.21$ فقط، لأن تغيّر \omterm{def:b1:geological-maps:attitude}{اتجاه الانحدار} على طبقة لطيفة يدير المستوي إدارة ضئيلة جدًا.

\textbf{12.} متوسط \omterm{def:b1:geological-maps:attitude}{الانحدار}: $1.68/\sqrt{12} = \ang{0.48}$؛ متوسط \omterm{def:b1:geological-maps:attitude}{اتجاه الانحدار}:
$2.94/\sqrt{12} = \ang{0.85}$.

\textbf{13.} لا: ف\ang{2.9} من \omterm{def:b1:geological-maps:attitude}{اتجاه الانحدار} على طبقة تنحدر بزاوية \ang{12} ليست إلا
$2.9\sin\ang{12.1} = \ang{0.6}$ من دوران المستوي، وهي أقلّ من \ang{1.7} من الانحدارات. وأي تموّج خفيف لسطح طبقي لطيف يبعثر اتجاه انحداره على نطاق واسع.

\textbf{14.} $D = \ang{-2.5}$: $83.1 - 2.5 = \ang{080.6}$؛ ومتوسطها هو
$080.6/12.1$، وهو متوسط الجيولوجي الأول نفسه.

\textbf{15.} يُقرأ \omterm{def:b1:geological-maps:attitude}{الانحدار} بمقياس الميل، على الشاقول الذي يعطيه الثقل؛ ولا تلعب الإبرة هنا أي دور.

\textbf{16.} نعم: فالدوران حول الشاقول يحرّك كل قراءة والمتوسط بالزاوية نفسها (\cref{exo:b1:field-lab-methods-1:11})؛ والجواب هو نفسه.

\textbf{17.} $10\times0.2 = \ang{2}$ من الخطأ في كل اتجاه انحدار، أكثر من ضعف الخطأ المعياري البالغ \ang{0.85}: خطأ منتظم لا يكشفه المتوسط.

\textbf{18.} $80.6 - 90 = \ang{350.6}$؛ \omterm{def:b1:geological-maps:attitude}{فالوضعية} هي امتداد \ang{351}، وانحدار \ang{12} نحو الشرق.

\textbf{19.} اتجاهات \omterm{def:b1:geological-maps:attitude}{الانحدار} $080.6$ و$080.6$: لا فرق بهذه الدقة؛ والانحدارات $12.1$ و$12.2$: \ang{0.1}.

\textbf{20.} النواظم $(\sin\ang{12.2}\sin\ang{80.6},
\sin\ang{12.2}\cos\ang{80.6}, \cos\ang{12.2})$ و$\vec S/\|\vec S\|$ تصنع زاوية قدرها \ang{0.13}.

\textbf{21.} $1.72/\sqrt{12} = \ang{0.50}$؛ والفرق، \ang{0.13}، هو ربع ذلك: غير دالّ. فالبوصلة والآبار متفقتان.

\textbf{22.} أربع مرات أكثر: 48 قراءة.

\textbf{23.} على سبيل المثال: ضبط انحراف خاطئ أو قديم؛ بوصلة منحرفة بالحديد القريب (مركبة، أدوات، صخور مغناطيسية)؛ مقياس ميل غير مضبوط؛ قراءات لم تُؤخذ إلا على أفضل الأسطح تعرّضًا، وقد لا تمثّل الجناح. ويُقلّل المتوسط الأخطاء العشوائية، لا المنتظمة.

\textbf{24.} \omterm{def:b1:field-lab-methods-1:mean-orientation}{الوضعية المتوسطة} للجناح الشرقي هي \textbf{081/12}: اتجاه انحدار \ang{081}، وانحدار \ang{12}.
\end{solution}
